\documentclass[12pt]{article}
\usepackage[T1]{fontenc}
\usepackage[utf8]{inputenc}
\usepackage{lmodern}
\usepackage{textcomp}
\usepackage{enumitem}
\usepackage{geometry}
\geometry{margin=1in}
\begin{document}
\begin{center}
Answer Key - Constitutional Law - Graduate Level Assessment 15 \
\small (Answers are ordered exactly as in the question paper.)
\end{center}

\section*{Multiple Choice}
\begin{enumerate}[label=\arabic*.]
\item \textbf{Answer:} B
\\ \textit{Options:} A) Law is unstable.; B) Law is determinate.; C) Law reflects economic power.; D) Law is politics.; E) Law is a social construct.
\\ \textit{Note:} Critical legal theorists argue that law is inherently indeterminate, meaning that legal outcomes are influenced by social, political, and economic factors rather than being fixed or certain.
\item \textbf{Answer:} A
\\ \textit{Options:} A) He claims that by giving priority to the needs of the poor, we can increase equality.; B) He rejects the idea of equality altogether.; C) He believes that equality is not a realistic goal.; D) He argues than an unequal society is inevitable.; E) He suggests that societal structures inherently support inequality.
\\ \textit{Note:} Parfit's opposition to strict equality stems from his belief that prioritizing the needs of the poor can lead to a more equitable society overall, as it addresses the systemic disadvantages faced by disadvantaged groups rather than merely equalizing resources regardless of context.
\item \textbf{Answer:} A
\\ \textit{Options:} A) Legal Pragmatism; B) Legal Formalism; C) Comparative; D) Analytical; E) Sociological
\\ \textit{Note:} Legal Pragmatism posits that the law should be understood as a tool for achieving practical outcomes and societal benefits, emphasizing the importance of what is considered "correct" in real-world contexts rather than strictly adhering to abstract principles or theories.
\item \textbf{Answer:} C
\\ \textit{Options:} A) Because the concept of punishment has evolved over time.; B) Because punishment can be justified in multiple ways.; C) Because any definition of punishment should be value-neutral.; D) Because the practice of punishment is separate from its justification.; E) Because the justification of punishment varies across cultures.
\\ \textit{Note:} Separating punishment from its justification is crucial because a value-neutral definition ensures that the concept can be analyzed and applied consistently within constitutional law, avoiding biases that may arise from moral or ethical considerations.
\item \textbf{Answer:} B
\\ \textit{Options:} A) Execution of the arrest warrant; B) investigative surveillance; C) Police interrogation prior to arrest; D) Post-charge lineups; E) Preliminary Hearing
\\ \textit{Note:} The correct answer is investigative surveillance because the Sixth Amendment right to counsel only attaches after formal charges have been initiated against an individual, and during investigative surveillance, no such charges have been filed.
\end{enumerate}

\section*{True / False}
\begin{enumerate}[label=\arabic*.]
\setcounter{enumi}{5}
\item \textbf{Answer:} False
\\ \textit{Options:} A) True; B) False
\\ \textit{Note:} Hume argued that natural law does not provide a sufficient foundation for state stability, as he believed that human conventions and practices, rather than abstract moral principles, are essential for maintaining order and security against external threats.
\item \textbf{Answer:} False
\\ \textit{Options:} A) True; B) False
\\ \textit{Note:} The protective principle of jurisdiction allows a state to assert jurisdiction over crimes committed outside its territory if they threaten the state's security or vital interests, thus making the statement false.
\item \textbf{Answer:} True
\\ \textit{Options:} A) True; B) False
\\ \textit{Note:} The international recognition of human rights after World War II, exemplified by documents like the Universal Declaration of Human Rights, reinvigorated natural law theories by emphasizing inherent human dignity and moral principles that transcend legal statutes, thus fostering a renewed focus on ethical foundations in law.
\item \textbf{Answer:} True
\\ \textit{Options:} A) True; B) False
\\ \textit{Note:} Partial performance of a contract does not terminate the contract by operation of law because, under contract law, parties are generally required to fulfill the terms of the agreement in full unless they mutually agree to terminate or modify the contract.
\item \textbf{Answer:} False
\\ \textit{Options:} A) True; B) False
\\ \textit{Note:} A two-thirds vote in the Senate, not a three-fourths vote, is required to convict and remove a president from office after impeachment by the House.
\end{enumerate}

\section*{Long Form Response}
\begin{enumerate}[label=\arabic*.]
\setcounter{enumi}{10}
\item \textbf{Answer:} In a trial concerning life insurance proceeds, a diary entry from the insured's mother documenting the insured's birth date may be deemed inadmissible due to hearsay rules. Hearsay is defined as an out-of-court statement offered to prove the truth of the matter asserted, which in this case is the accuracy of the birth date. Since the mother is not available to testify and the entry lacks corroborating evidence, it fails to meet the reliability standards necessary for admissibility. Furthermore, the diary entry may not fall under any established exceptions to the hearsay rule, such as present sense impression or excited utterance, thereby reinforcing its exclusion from the evidentiary record. This analysis underscores the importance of ensuring that evidence presented in court adheres to procedural rules that protect the integrity of the judicial process.
\\ \textit{Note:} correct because the diary entry is considered hearsay, as it is an out-of-court statement used to establish the truth of the insured's birth date without the mother's testimony or corroborating evidence, thus failing to meet the necessary reliability standards for admissibility in court.
\item \textbf{Answer:} Legal positivism asserts that the relationship between law and morality is complex, often emphasizing that while laws may reflect moral values, they are ultimately distinct entities. The claim that legal positivism regards morals and law as inseparable highlights the view that legal systems can incorporate moral standards, yet the validity of law does not inherently depend on moral correctness. This perspective has significant implications for constitutional law, as it suggests that constitutional provisions may be interpreted within a moral framework, yet their enforcement and legitimacy remain grounded in their adherence to established legal processes rather than moral considerations alone. Consequently, legal positivism fosters a critical examination of how moral values influence constitutional interpretation while maintaining a clear boundary between law and ethics.
\\ \textit{Note:} The answer correctly identifies that legal positivism differentiates between law and morality, emphasizing that while laws may reflect moral values, their validity does not depend on moral correctness, which has important implications for how constitutional law can incorporate or reflect moral standards without being bound to them.
\end{enumerate}

\vfill
\noindent\textit{End of Answer Key}
\end{document}